% 本书共享导言区；第 1 章以原书实际编号为准。
\usepackage{amsmath,amssymb,amsthm,mathtools,bm}
\usepackage{tikz}
\usetikzlibrary{arrows.meta,positioning,calc,backgrounds,shapes.geometric}
\usepackage{graphicx,booktabs,array,tabularx,enumitem,xcolor,caption,fancyhdr,placeins,needspace}
\usepackage[a4paper,top=2.6cm,bottom=2.6cm,left=2.6cm,right=2.6cm,headsep=0.7cm,footskip=1.0cm]{geometry}
\setmainfont{TeX Gyre Termes}
\definecolor{inkblue}{RGB}{28,62,110}
\definecolor{plotgray}{RGB}{35,31,32}
\linespread{1.32}
\setlength{\parskip}{0.25em}
\setlength{\emergencystretch}{1.2em}
\tolerance=2000
\setlength{\headheight}{15pt}
\ctexset{
 chapter={format=\centering\heiti\zihao{2},name={第,章},number=\chinese{chapter},beforeskip=1.2cm,afterskip=1.2cm},
 section={format=\heiti\zihao{3},beforeskip=1.1ex plus .3ex minus .2ex,afterskip=.8ex plus .2ex},
 subsection={format=\heiti\zihao{4},beforeskip=1.0ex plus .3ex minus .2ex,afterskip=.6ex plus .2ex},
 figurename=图,tablename=表
}
\captionsetup{font=small,labelsep=quad,labelfont=bf,skip=8pt}
\captionsetup[figure]{justification=justified,singlelinecheck=false}
\captionsetup[table]{justification=justified,singlelinecheck=false}
\pagestyle{fancy}
\fancyhf{}
\fancyhead[LE]{\small\kaishu\leftmark}
\fancyhead[RO]{\small\kaishu\rightmark}
\fancyfoot[C]{\small\thepage}
\renewcommand{\headrulewidth}{0.4pt}
\fancypagestyle{plain}{\fancyhf{}\fancyfoot[C]{\small\thepage}\renewcommand{\headrulewidth}{0pt}}
\makeatletter
\renewcommand{\chaptermark}[1]{\markboth{\CTEXthechapter\quad #1}{}}
\renewcommand{\sectionmark}[1]{\markright{\thesection\quad #1}}
\makeatother
% 原书各类结果按各自的“章.序号”编号；逐条保留源编号时用 bookresult。
\newtheoremstyle{cnresult}{7pt}{7pt}{\normalfont}{}{\bfseries}{}{0.7em}{}
\theoremstyle{cnresult}
\newtheorem{definition}{定义}[chapter]
\newtheorem{theorem}{定理}[chapter]
\newtheorem{lemma}{引理}[chapter]
\newtheorem{corollary}{推论}[chapter]
\newtheorem{proposition}{命题}[chapter]
\newtheorem{remark}{注}[chapter]
\newtheorem{example}{例}[chapter]
\newtheorem{assumption}{假设}[chapter]
\newtheorem{conjecture}{猜想}[chapter]
\newtheorem{algorithm}{算法}[chapter]
\newtheorem{problem}{习题}[chapter]
\newenvironment{bookresult}[2]{\par\addvspace{7pt}\noindent\textbf{#1 #2.}\enspace\ignorespaces}{\par\addvspace{7pt}}
\newcommand{\chapterauthors}[1]{\begin{center}\normalsize #1\end{center}\vspace{.6em}}
\newcommand{\chapteraffiliations}[1]{\begingroup\footnotesize\noindent #1\par\endgroup\vspace{1em}}
\renewcommand{\proofname}{证明}
\newenvironment{remarkenv}{\par\noindent\textbf{注。}\ignorespaces}{\par}
\newcommand{\PA}{\mathrm{PA}}
\newcommand{\pa}{\mathrm{pa}}
\newcommand{\ND}{\mathrm{ND}}
\newcommand{\AN}{\mathrm{AN}}
\newcommand{\DE}{\mathrm{DE}}
\newcommand{\EE}{\mathbb{E}}
\newcommand{\RR}{\mathbb{R}}
\DeclareMathOperator{\Var}{Var}
\DeclareMathOperator{\Cov}{Cov}
\DeclareMathOperator{\tr}{tr}
\DeclareMathOperator{\diag}{diag}
\DeclareMathOperator{\supp}{supp}
\DeclareMathOperator{\sign}{sign}
\usepackage{listings}
\lstset{basicstyle=\ttfamily\footnotesize,columns=fullflexible,keepspaces=true,
 breaklines=true,numbers=left,numberstyle=\tiny\color{gray},numbersep=6pt,
 frame=single,framerule=.3pt,rulecolor=\color{gray!50},showstringspaces=false,
 tabsize=2,aboveskip=6pt,belowskip=6pt}
\usepackage[most]{tcolorbox}
\newtcolorbox{sourcebox}[1][]{enhanced,breakable,colback=gray!3,colframe=gray!50,
 boxrule=.4pt,arc=0pt,left=7pt,right=7pt,top=5pt,bottom=5pt,title={#1},coltitle=black,
 colbacktitle=gray!10,fonttitle=\bfseries}
\usetikzlibrary{decorations.pathreplacing,decorations.markings,fit,patterns}
\tikzset{
 vtxb/.style={vtx,very thick},
 vtxs/.style={vtx,minimum size=5.5mm,font=\footnotesize},
 lat/.style={latent},
 edb/.style={->,very thick},
 bid/.style={<->,semithick},
 bidd/.style={<->,semithick,dashed}
}

\newcommand{\indep}{\mathrel{\perp\!\!\!\perp}}
\newcommand{\nindep}{\mathrel{\not\!\perp\!\!\!\perp}}
\DeclareMathOperator*{\argmax}{arg\,max}
\DeclareMathOperator{\doop}{do}
\newcommand{\dox}[1]{\mathrm{do}(#1)}
\def\doxp(#1){\mathrm{do}(#1)}
\tikzset{
 causal/.style={>={Stealth[length=2.1mm,width=1.5mm]},semithick},
 vtx/.style={circle,draw,inner sep=1pt,minimum size=7mm,font=\small},
 latent/.style={vtx,dashed},
 ed/.style={->,semithick},
 edd/.style={->,semithick,dashed}
}
\setlist[enumerate]{leftmargin=2.3em,itemsep=0.35em,topsep=0.4em}
\renewcommand{\arraystretch}{1.3}
\setcounter{secnumdepth}{2}
\setcounter{tocdepth}{2}
\usepackage{longtable,multicol}
\usepackage{url}
\usepackage[bookmarks=true,bookmarksnumbered=true,unicode=true,colorlinks=true,linkcolor=inkblue,citecolor=inkblue,urlcolor=inkblue]{hyperref}
\hypersetup{
 pdftitle={数据科学中的靶向学习：复杂纵向研究的因果推断（中文译稿）},
 pdfauthor={Mark J. van der Laan; Sherri Rose},
 pdfsubject={Targeted Learning in Data Science (2018), Chinese translation},
 pdfkeywords={靶向学习,TMLE,纵向研究,因果推断}
}

\usepackage{pdflscape,adjustbox,rotating}
\DeclareMathOperator*{\argmin}{arg\,min}
\DeclareMathOperator{\expit}{expit}
\DeclareMathOperator{\logit}{logit}
\DeclareMathOperator{\E}{E}
\DeclareMathOperator{\bias}{Bias}
\DeclareMathOperator{\MSE}{MSE}
\newcommand{\1}{\mathbb{I}}
\newcommand{\Pn}{P_n}
\newcommand{\Pzero}{P_0}
\newcommand{\Gbar}{\overline{G}}
\newcommand{\Qbar}{\overline{Q}}
\newcommand{\Lbar}{\overline{L}}
\newcommand{\Abar}{\overline{A}}
\setlength{\LTpre}{6pt}
\setlength{\LTpost}{6pt}
